\abstractpageen{Abstract}{%
  This thesis addresses a constrained, instrument-family-level version of
  automatic musical instrument recognition from audio recordings using modern
  machine learning methods. The developed system is capable of classifying
  seven audio classes in real time (guitar, piano, vocals, strings, reeds,
  brass, and ambient noise/silence) from monophonic or dominant-source audio
  signals.

  The proposed approach relies on the extraction of three types of spectral
  features (Log-Mel spectrograms, STFT, and MFCC), which are fed into
  two-dimensional convolutional neural networks (2D~CNN) trained in the
  Google Colab environment. The dataset consists of self-collected
  internet-sourced instrument clips, noise/silence samples from ESC-50,
  software-augmented training samples, and self-recorded microphone playback
  samples. A data augmentation module approximates real-world recording
  conditions by adding reverberation, ambient noise, and microphone
  characteristics.

  Although the experimental results are far from perfect, they show promising
  accuracy on the test data considering the small, self-assembled dataset, and
  the system can also be run as a real-time microphone-based demonstration. The
  work demonstrates the feasibility of building a measurable and reproducible
  instrument-family-level audio classification prototype using accessible
  computational resources.
}{%
\textit{\textbf{Keywords:}} musical instrument recognition, convolutional neural
network, Mel-spectrogram, STFT, MFCC, audio classification, waveform, spectral image%
}
